% 对应题号 M1-17；图与本题条件、推导、结论逐项对应。
% 图契约：印刷宽16cm；最小字号8.4pt；几何/场或曲线在左，计算关系在右。
\documentclass[tikz,border=7pt]{standalone}
\usepackage{fontspec,xeCJK,amsmath,amssymb,graphicx}
\IfFontExistsTF{Times New Roman}{\setmainfont{Times New Roman}}{\setmainfont{TeX Gyre Termes}}
\IfFontExistsTF{Heiti SC}{\setCJKmainfont[AutoFakeBold=1.5]{Heiti SC}}{\setCJKmainfont[AutoFakeBold=1.5]{FandolHei-Regular}}
\usetikzlibrary{arrows.meta,bending,calc,patterns,positioning}
\definecolor{ink}{HTML}{243745}
\definecolor{blue}{HTML}{275DAB}
\definecolor{green}{HTML}{177C63}
\definecolor{red}{HTML}{BD4935}
\definecolor{orange}{HTML}{AA6C16}
\definecolor{purple}{HTML}{715393}
\tikzset{every node/.style={font=\fontsize{9.3}{12.5}\selectfont,text=ink},
 axis/.style={-{Stealth[length=4pt]},line width=.65pt,draw=ink!65},
 vec/.style={-{Stealth[length=5pt]},line width=1pt,draw=blue},
 short/.style={-{Stealth[length=3pt]},line width=.8pt},
 guide/.style={dashed,draw=ink!40,line width=.5pt},
 card/.style={rounded corners=2pt,draw=blue!65,fill=blue!4,align=center,text width=6.6cm,inner xsep=9pt,inner ysep=7pt},
 title/.style={font=\fontsize{10.5}{14}\selectfont\bfseries,text width=15cm,align=center},
 label/.style={font=\fontsize{8.4}{11}\selectfont,align=center}}
\begin{document}
\begin{tikzpicture}[x=1cm,y=1cm]
\path[use as bounding box] (-8,-3.4) rectangle(8,3.4);
\node[title] at (0,3) {M1-17\quad 小位移与梯度的有向投影};
\begin{scope}[shift={(-4,0)}]
\draw[axis] (-2.4,0)--(2.5,0) node[right] {$x$};
\draw[axis] (0,-2)--(0,2.2) node[above] {$y$};

\foreach \r in {1,1.7,2.3}{\draw[blue!35] (0,0) circle(\r);}
\coordinate(P) at (.7,1.4);\fill[red] (P) circle(2pt) node[below left=1pt] {$P$};
\draw[vec] (P)--(1.15,2.3) node[right] {$\nabla f(P)$};
\draw[vec,orange] (P)--(1.9,.2) node[right] {$\Delta\mathbf r$};
\node[label] at (0,-2.7) {位移箭头放大示意，实际为 $(.01,-.01)$};

\end{scope}\node[card,draw=blue!65,fill=blue!4] at (3.95,1.65) {$f=x^2+y^2$，$\nabla f(1,2)=(2,4)$};
\node[card,draw=green!65,fill=green!4] at (3.95,0) {$\Delta f\simeq(2,4)\cdot(.01,-.01)=-.02$};
\node[card,draw=orange!65,fill=orange!4] at (3.95,-1.65) {位移含向内的法向分量，\\函数值在一阶近似下减小。};
\end{tikzpicture}\end{document}
